\chapter{\'Etat de l'art et choix technologiques}
\finPageTitre

\section{Introduction}
La r\'eussite d'un projet d'ing\'enierie logicielle distribu\'ee repose sur l'adoption d'une d\'emarche m\'ethodologique rigoureuse et sur la s\'election d'un \'ecosyst\`eme technologique adapt\'e aux exigences de performance et de scalabilit\'e. Ce deuxi\`eme chapitre pr\'esente l'approche m\'ethodologique retenue, d\'eveloppe une \'etude comparative approfondie des paradigmes de communication et des courtiers de messages, et formalise l'architecture microservices du syst\`eme d'ingestion IoT.

\section{Approche m\'ethodologique et organisation du projet}
\subsection{Cycle de d\'eveloppement it\'eratif et agile}
Afin d'assurer une convergence rapide vers les objectifs op\'erationnels, une m\'ethodologie de d\'eveloppement agile inspir\'ee du cadre Scrum a \'et\'e appliqu\'ee. Cette approche s'est concr\'etis\'ee par des cycles it\'eratifs bi-hebdomadaires permettant d'int\'egrer en continu les retours des encadrants technique et acad\'emique.

\subsection{Principes d'ing\'enierie logicielle appliqu\'es}
Le d\'eveloppement s'est conform\'e aux principes de conception logicielle reconnus :
\begin{itemize}
    \item \textbf{S\'eparation des pr\'eoccupations} : D\'ecouplage strict entre la r\'eception des trames, la validation m\'etier et la persistance ;
    \item \textbf{Autonomie des microservices} : Chaque service dispose de sa propre base de donn\'ees ou de son espace logique d\'edi\'e selon le patron \textit{Database-per-Service}~\cite{richardson2018microservices} ;
    \item \textbf{R\'esilience et tol\'erance aux pannes} : Mise en place de circuits d'interruption et de m\'ecanismes de relance automatique.
\end{itemize}

\section{\'Etude comparative des technologies de communication et de courtage}
\subsection{\'Evaluation des protocoles de transport IoT}
La transmission des donn\'ees depuis les passerelles p\'eriph\'eriques vers la plateforme centrale a n\'ecessit\'e la comparaison entre deux protocoles majeurs :
\begin{itemize}
    \item \textbf{HTTP/REST} : Protocole synchrone, reposant sur TCP. Tr\`es r\'epandu et ais\'e \`a d\'eboguer, il induit n\'eanmoins une surcharge d'en-t\^etes cons\'equente et une d\'ependance directe \`a la disponibilit\'e imm\'ediate du r\'ecepteur ;
    \item \textbf{MQTT} : Protocole l\'eger de publication/abonnement con\c{c}u sp\'ecifiquement pour les environnements \`a bande passante contrainte. Sa faible empreinte r\'eseau et ses niveaux de qualit\'e de service (QoS) en font un candidat optimal pour la t\'el\'em\'etrie capteurs.
\end{itemize}

\subsection{S\'election du courtier de messages distribu\'e}
Pour absorber les pics de trafic sans impacter les services aval, une couche de courtage asynchrone est indispensable. Le Tableau~\ref{tab:synthese-technologies} synth\'etise la comparaison \'etablie entre les trois solutions \'etudi\'ees.

\begin{table}[H]
\centering
\begin{tabularx}{\textwidth}{p{3.2cm} X X p{3.2cm}}
\toprule
\textbf{Courtier} & \textbf{Avantages Principaux} & \textbf{Limites Notables} & \textbf{Choix Retenu} \\
\midrule
\textbf{RabbitMQ} & Routage flexible, support AMQP \'etendu & Moins performant sur les flux tr\`es massifs & Non retenu \\
\textbf{Redis Streams} & Vitesse en m\'emoire, latence tr\`es basse & Persistance secondaire sur disque & Utilis\'e en cache \\
\textbf{Apache Kafka} & D\'ebit exceptionnel, tol\'erance aux pannes, r\'etention persistante & Gestion plus complexe de cluster & \textbf{Retenu pour le c\oe ur d'ingestion} \\
\bottomrule
\end{tabularx}
\caption{Comparaison technique des solutions de courtage de messages}
\label{tab:synthese-technologies}
\end{table}

L'analyse comparative consign\'ee dans le Tableau~\ref{tab:synthese-technologies} montre qu'Apache Kafka offre les garanties d'immuabilit\'e et de partitionnement indispensables au traitement parall\`ele de millions de mesures journali\`eres~\cite{newman2015building}.

\subsection{Moteur de persistance pour s\'eries temporelles}
Les flux IoT g\'en\`erent des s\'eries chronologiques n\'ecessitant des fonctions analytiques avanc\'ees (agr\'egations temporelles, fen\^etres glissantes). Le choix s'est port\'e sur \textbf{TimescaleDB}, extension sp\'ecialis\'ee pour PostgreSQL, alliant l'expressivit\'e du mod\`ele relationnel et l'indexation optimis\'ee des s\'eries temporelles via des \textit{hypertables}.

\section{Architecture globale et mod\'elisation microservices}
\subsection{Patron d'architecture r\'eactive orient\'ee \'ev\'enements}
L'architecture globale du syst\`eme repose sur une topologie orient\'ee \'ev\'enements (\textit{Event-Driven Architecture}). La Figure~\ref{fig:architecture-globale} sch\'ematise l'agencement des microservices et la circulation des messages.

\begin{figure}[H]
\centering
\IfFileExists{figures/ch2/architecture_globale.png}{%
    \includegraphics[width=0.85\textwidth]{figures/ch2/architecture_globale.png}%
}{%
    \fbox{\parbox[c][5cm][c]{0.85\textwidth}{\centering \textit{[Sch\'ema de l'architecture microservices : Ingestion MQTT/HTTP $\rightarrow$ Kafka $\rightarrow$ Processeurs de calcul $\rightarrow$ TimescaleDB]}}}%
}
\caption{Architecture microservices distribu\'ee pour l'ingestion IoT}
\label{fig:architecture-globale}
\end{figure}

Comme l'illustre la Figure~\ref{fig:architecture-globale}, la passerelle frontale re\c{c}oit les paquets et les \'emet imm\'ediatement dans une file d'attente tamponn\'ee, garantissant une r\'eponse quasi-instantan\'ee au client tout en d\'ecouplant compl\`etement les modules d'enrichissement et d'indexation.

\subsection{D\'ecomposition des responsabilit\'es m\'etier}
Le d\'ecoupage fonctionnel aboutit \`a trois microservices autonomes :
\begin{itemize}
    \item \textbf{Gateway-Service} : Service frontal charg\'e de l'authentification des capteurs et de la conversion des trames brutes en \'ev\'enements normalis\'es ;
    \item \textbf{Ingestion-Processor} : Composant d'enrichissement m\'etier calculant les moyennes mobiles et d\'etectant les anomalies de d\'epassement de seuil ;
    \item \textbf{Storage-Sink} : Service assurant l'\'ecriture par lots (\textit{batching}) dans TimescaleDB pour r\'eduire la contention I/O.
\end{itemize}

\section{D\'emarche d'assurance qualit\'e logicielle}
L'ensemble des modules a \'et\'e soumis \`a une d\'emarche qualit\'e continue :
\begin{itemize}
    \item Tests unitaires automatis\'es couvrant la logique de d\'ecodage des trames ;
    \item Conteneurisation int\'egrale via des images Docker optimis\'ees \`a \'etapes multiples (\textit{multi-stage builds}) ;
    \item Respect des r\`egles de nommage et absence totale de code redondant selon le principe DRY (\textit{Don't Repeat Yourself}).
\end{itemize}

\sectionTransition{%
Ce deuxi\`eme chapitre a justifi\'e les choix technologiques et pr\'esent\'e la mod\'elisation de l'architecture microservices. Le troisi\`eme chapitre expose l'impl\'ementation pratique des composants, l'environnement de d\'eploiement ainsi que l'\'evaluation exp\'erimentale des performances sous charge.
}
